\documentclass[10pt,a4paper]{amsart}
\usepackage[margin=19mm]{geometry}
\usepackage{amsmath,amssymb,mathtools}
\usepackage[scr=rsfs,cal=euler]{mathalfa}
\usepackage{xfrac}
\usepackage[unicode,hidelinks,bookmarks=false]{hyperref}
\newcommand{\ie}{i.e.\ }
\newcommand{\cf}{cf.\ }
\newcommand{\resp}{respectively\ }
\newcommand{\Subsection}{\subsection}
\providecommand{\nobreakdash}{\nobreak\hbox{-}}
\setlength{\emergencystretch}{2em}
\multlinegap0pt
\usepackage{bm}
\usepackage{amssymb}
\usepackage{float}
\newtheorem{proposition}{Proposition}
\usepackage{cleveref}
\crefname{proposition}{Proposition}{Propositions}
\def\C{{\mathbb C}}
\def\CZ{\rm CZ}
\def\N{{\mathbb N}}
\def\Q{{\mathbb Q}}
\DeclareMathOperator{\SOB}{SOB}
\def\rd{{\rm d}}
\DeclareMathOperator{\symp}{Symp}
\title{Algebraic planar torsion in contact manifolds}
\author{Zhengyi Zhou}
\date{Source: arXiv:2603.06031v1 (CC BY 4.0)}
\begin{document}
\maketitle

\noindent\emph{Source and licence.} Algebraic planar torsion in contact manifolds. Version arXiv:2603.06031v1 (CC BY 4.0), distributed by arXiv under the Creative Commons Attribution 4.0 International licence, http://creativecommons.org/licenses/by/4.0/, as recorded in the arXiv licence metadata for this exact version and shown on the abstract page https://arxiv.org/abs/2603.06031v1. Full licence notice: \texttt{licenses/arxiv-2603.06031-CC-BY-4.0.txt}.\par\smallskip

\noindent\textit{Editorial scope.} Counted unit: the article's proposition prop:filtration (source0.tex lines 1938--1940). The article's complete original proof of it is delivered with it (source0.tex lines 1941--1952, one \texttt{proof} environment of 12 lines). No mathematics is written, repaired or re-derived: the statement and the proof are the article's own lines, in the article's own notation, and no retained line is substituted or re-flowed.  Retained uncounted as the source-local context the proof needs: lines 696--698; lines 1902--1904; lines 1915--1922.  Nothing was omitted from any retained span, so the package carries no omission record.  No retained line contains a citation, so the wrapper carries no bibliography and no bibliography entry.  No retained line contains a figure environment, an image-inclusion command or drawing code, so no image is delivered and the package declares no asset dependency.  Editorial formatting only, in addition: an AMS \texttt{amsart} preamble that restates the article's own declarations and macros used by the retained lines, namely the environments \texttt{proposition} on separate counters, together with the standard packages those lines need.  The retained environments print in this document as Proposition 1 in order of appearance, whereas the article prints the same items with the numbering of its own full document.  The complete native source of the article as distributed in the arXiv e-print for this exact version is bundled unchanged as \texttt{sources/195837/source0.tex}.
\begin{equation}\label{eqn:TW}
    \alpha = K\lambda_S+\alpha_M \text{ on } S_{k}\times M, \quad \alpha = K\rd t + \lambda_{i}+\beta(t)\eta_{i} \text{ on } (\Sigma_i)_{\phi_i}.
\end{equation}

\begin{proposition}\label{prop:c1}
    Assume that there exists a trivialization $\Psi: \det_{\C}\oplus^N T\Sigma\to \underline{\C}$ for $N\ge 1 \in \N$, which induces a trivialization $\Psi|_{\partial \Sigma}$ of $\det_{\C}\oplus^N \xi_{\partial \Sigma}$ on $\partial \Sigma$ (where $\xi_{\partial \Sigma}$ is the contact structure on $\partial \Sigma$), such that the trivialization $(\psi_i^{-1}\circ \psi_1)_*\Psi|_{\partial \Sigma}$ can be extended to $\Sigma$ for all $1\le i \le k=|\bm{\phi}|$. Then we have $c_1^{\Q}(\SOB(\Sigma,\bm{\phi},\bm{\psi}))=0$.
\end{proposition}

\begin{proposition}\label{prop:CZ}
    Let $\alpha_M$ be a contact form on $M$ and let $D>0$ be a number not in the Reeb spectrum of $\alpha_M$ (the set of periods of Reeb orbits of $\alpha_M$). There exists a contact form on $\SOB(\Sigma,\bm{\phi},\bm{\psi})$, which is a slight perturbation and smoothing of \eqref{eqn:TW}, such that the Reeb orbits of period $<D$ are either
    \begin{enumerate}
        \item contained in the paper or part of the spine over the collar neighborhoods of $\partial S_k$, which project to the $m$-fold cover of a boundary component of $S_k$ through $\pi_S$, for $m>0$;
        \item or contained in the spine outside of the collar neighborhoods above, in the form of $\hat{\gamma},\check{\gamma}_1,\ldots,\check{\gamma}_{k-1}$ for each Reeb orbit $\gamma$ of $M$ with period $<D$, where $k=|\bm{\phi}|$. Those orbits are copies of $\gamma \subset M$ over a point in $S_k$.
    \end{enumerate}
    Assume the conditions in \Cref{prop:c1} hold. Then we have $\mu_{\CZ}(\hat{\gamma})=\mu_{\CZ}(\gamma)+1$ and $\mu_{\CZ}(\check{\gamma}_i)=\mu_{\CZ}(\gamma)$, where $\mu_{\CZ}(\gamma)$ is the Conley-Zehnder index computed using $\Psi|_{\partial \Sigma}$ (which are rational in general).
\end{proposition}

\begin{proposition}\label{prop:filtration}
    Let $\alpha^1_M>\alpha^2_M>\ldots >\alpha^\infty_M$ be a sequence of contact forms on $M$. Then the construction in \Cref{prop:CZ} yields contact forms $\alpha_1>\alpha_2>\ldots$ on $\SOB(\Sigma,\bm{\phi},\bm{\psi})$, with period threshold $D_i$ converging to $\infty$.
\end{proposition}

\begin{proof}
    Without loss of generality, we may assume $\phi_i\in \symp_c(\Sigma_i,\rd \lambda_i)$, where the Liouville form $\lambda_i$ restricted to $\partial \Sigma_i$ is $\psi_i^*\alpha^{\infty}_M$. We fix $K\gg 0$ such that the following holds:
    \begin{enumerate}
        \item $K\rd t + \beta(t)\eta_i+\lambda_i$ is a contact form on $\Sigma_{\phi_i}$;
        \item $K>\beta'(t)\eta(X_i(t))$ for all $t\in S^1$, where $X_i(t)\in T\Sigma$ is defined by $\iota_{X_i(t)}\rd \lambda_i=\lambda_i+\beta(t)\eta_i$ with $\phi_i^*\lambda_i=\lambda_i+\eta_i$.
    \end{enumerate}
    Consider the following domain:
    $$\left(\widehat{S}_k \times M\times [1,\infty)_r \bigcup_{i=1}^k [1,\infty)_p \times \Sigma_{\phi_i}, \quad \lambda:=\left(\widehat{\lambda}_S+r\alpha_M \bigcup_{i=1}^k Kp\rd t + \beta(t)\eta_i+\lambda_i\right)\right).$$
    With the conditions above, $\lambda$ is a Liouville form, and the Liouville vector field $X_\lambda$ is defined by 
    $$\left(X_S+r\partial_r\right) \bigcup_{i=1}^k \left((Kp-\beta'(t)\eta(X_i(t)))\partial_p+X_i(t)\right),$$
    where $X_S$ is the Liouville vector field on the Liouville completion $\widehat{S}_k$. The second condition guarantees that $X_{\lambda}$ has no zeros on the domain above. Note that the contact forms in \Cref{prop:CZ} are the restrictions of $\lambda$ to a smoothing and perturbation of the hypersurface given by $r=h_i>h_{i+1}>1$ for $\alpha^i_M=h_i\alpha^{\infty}_M$ and $p=1+\delta_i$ with $\delta_1>\delta_2>\ldots >0$. The Liouville vector field $X_\lambda$ is transverse to those hypersurfaces, making the region between them a Liouville cobordism, homotopic to the trivial cobordism. Consequently, the restriction of $\lambda$, namely $\alpha_i$, satisfies $\alpha_i>\alpha_{i+1}$ under a suitable identification of the hypersurfaces (i.e., using the flow of $X_{\lambda}$), such that the period threshold $D_i$ converges to $\infty$ by flattening the function $f$ in the perturbations.
\end{proof}

\end{document}
